% !TEX TS-program = lualatex
% !TEX encoding = UTF-8

\documentclass[VesperaleOP.tex]{subfiles}

\ifcsname preamble@file\endcsname
  \setcounter{page}{\getpagerefnumber{M-vop50_sanctorale_decembris}}
\fi

\begin{document}

\festum{}{Proprium Sanctorum}{}{Proprium Sanctorum}{0}
\addcontentsline{toc}{chapter}{Proprium Sanctorum}

\festum{die 30 novembris}{S. Andreæ Apostoli}{II classis}{die 30 novembris --- S. Andreæ Apostoli}{2}
\titulumrubrica{Ad Vesperas}
\cantus[Ant. I b]{1130VA}{Super Ps.}
\rubrica{Psalmi, capitulum, hymnus et \vvrub de Communi Apostolorum, pag. XXX}
\cantus{1130VAM}{Ad Magn.}
\oratio{Maiestátem tuam, Dómine, supplíciter exorámus: ut, sicut Ecclésiæ tuæ beátus Andréas Apóstolus éxstitit prædicátor et rector,\pscross{}ita apud te sit pro nobis perpétuus intercéssor.\psstar{} Per.}


\festum{2.}{S. Bibianæ Virginis et Martyris}{III classis}{die 2 decembris --- S. Bibianæ}{3}
\oratio{Indulgéntiam nobis, quǽsumus, Dómine, beáta Bibiána Virgo et Martyr implóret:\psstar{} quæ tibi grata semper éxstitit,\pscross{}et mérito castitátis et tuæ professióne virtútis.\psstar{} Per.}


\festum{3.}{S. Francisci Xaverii Confessoris}{III classis}{die 3 decembris --- S. Francisci Xaverii}{3}
\oratio{Deus, qui Indiárum gentes beáti Francísci prædicatióne et miráculis Ecclésiæ tuæ aggregáre voluísti:\psstar{} concéde propítius ut, cuius gloriósa mérita venerámur,\pscross{}virtútum quoque imitémur exémpla.\psstar{} Per.}


\festum{4.}{S. Petri Chrysologi Confessoris}{III classis}{die 4 decembris --- S. Petri Chrysologi}{3}
\oratio{Deus, qui beátum Petrum Chrysólogum Doctórem egrégium, divínitus præmonstrátum, ad regéndam et instruéndam Ecclésiam tuam éligi voluísti:\psstar{} præsta, quǽsumus, ut, quem Doctórem vitæ habúimus in terris,\pscross{}intercessórem habére mereámur in cælis.\psstar{} Per.}


\festum{6.}{S. Nicolai Episcopi et Confessoris}{III classis}{die 6 decembris --- S. Nicolai}{3}
\titulumrubrica{Ad Vesperas}
\cantus[Ant. I b]{1206VA}{Super Ps.}
\rubrica{Psalmi de Communi unius Martyris, pag. XXX}
\rubrica{Capitulum, hymnus et \vvrub de Communi unius Confessoris, pag. XXX}
\cantus{1206VAM}{Ad Magn.}
\oratio{Deus, qui beátum Nicoláum, Pontíficem tuum, innúmeris decorásti miráculis;\pscross{}tríbue, quǽsumus, ut eius méritis et précibus a gehénnæ incéndiis liberémur.\psstar{} Per.}

%%%%DIE 8 DEC. IN CONCEPTIONE IMM. B.M.V.
\festum{die 8 decembris}{In Conceptione Immaculata\\Beatæ Mariæ Virginis}{I classis}{die 8 decembris --- In Conceptione Imm. B. M. V.}{1}
\titulumrubrica{Ad I Vesperas}
\cantus[Ant. VII b]{1207VA}{Super Ps.}
\capitulum{Dóminus possédit me in inítio viárum suárum, ántequam quidquam fáceret a princípio.\pscross{}Ab ætérno ordináta sum et ex antíquis, ántequam terra fíeret.\psstar{} Nondum erant abýssi, et ego iam concépta eram.\pscross{}}
\cantus{1207VR}{\rrann}
\rubrica{Hic genuflectitur usque ad \normaltext{Sumens illud Ave.}}
\cantus{CBVMVH}{Hymnus}

\versiculus{Immaculáta Concéptio est hódie sanctæ Maríæ Vírginis.}{Quæ serpéntis caput virgíneo pede contrívit.}

\cantus{1207VAM}{Ad Magn.}
\oratio{Deus, qui per immaculátam Vírginis Conceptiónem dignum Fílio tuo habitáculum præparásti:\psstar{} quǽsumus ut, qui ex morte eiúsdem Fílii tui prævísa eam ab omni labe præservásti,\pscross{}nos quoque mundos eius intercessióne ad te perveníre concédas.\psstar{} Per.}
\rubrica{Et fit commemoratio feriæ.}

\titulumrubrica{Ad II Vesperas}
\cantus[Ant. VII b]{1208VA}{Super Ps.}
\rubrica{Psalmi ut in festis B. Mariæ V. p. XXX.}

\rubrica{Capitulum, hymnus et \vvrub ut in I Vesperis.}

\cantus{1208VAM}{Ad Magn.}

\rubrica{Oratio ut in I Vesperis et fit commemoratio feriæ tantum.}
\lineamagna
%%%%



\end{document}